\section{解决方案}

\subsection{总体架构设计}
架构以\textbf{数据资产可复现、证据可追溯、决策可复核}为主线。系统分为数据治理与分析资产层、模型与索引层、服务与接口层、智能体编排层四类能力;各层既可独立评估,也共享同一数据口径,便于持续更新与复现。
图~\ref{fig:project-architecture} 展示底层数据到智能体交互的完整链路,图~\ref{fig:capability-map} 则把工程组件收束为评委可快速识别的能力地图。

\begin{figure}[!htbp]
\centering
\begin{tikzpicture}[
  x=1mm, y=1mm, font=\small,
  layer/.style={draw=SciLine, rounded corners=1.4pt, line width=0.6pt, align=center, inner sep=4pt, minimum height=11mm, text=SciInk, fill=white},
  band/.style={draw=SciLine, fill=SciPanel, rounded corners=1.4pt, line width=0.45pt, align=center, inner sep=3pt, minimum height=10mm, text=SciInk},
  arrow/.style={-{Latex[length=2mm]}, line width=0.7pt, draw=SciMuted},
  lab/.style={font=\scriptsize\bfseries, text=SciTealDark, anchor=east}
]
% layers bottom-up
\node[lab] at (-3,5)   {1. 采集};
\node[layer, text width=150mm] (raw) at (78,5) {6 公开来源(arXiv/PubMed/PMC/OpenAlex/Crossref/DOAJ)→ 原始 JSONL(\texttt{data/raw\_canonical})};

\node[lab] at (-3,22)  {2. 治理};
\node[layer, text width=150mm] (proc) at (78,22) {清洗规范化 · \textbf{标题优先去重} · front-matter/PDF 垃圾过滤 · NUL 清洗 → 159{,}164 篇处理后语料(\texttt{processed})};

\node[lab] at (-3,42)  {3. 加工};
\node[band, text width=70mm] (ana) at (40,42) {\textbf{分析层}\\领域分布·关键词演化·作者网络·主题模型\\{\scriptsize ->《数据分析报告》(成果 1)}};
\node[band, text width=72mm] (model) at (117,42) {\textbf{模型层 / RAG}\\片段切分·嵌入向量·趋势·推荐·图谱\\{\scriptsize -> 科研智能体"模型文件"(成果 2)}};

\node[lab] at (-3,60)  {4. 服务};
\node[layer, text width=150mm] (pg) at (78,60) {PostgreSQL + pgvector 服务层:papers/chunks/edges · 全文索引(tsvector)· 向量索引(ivfflat)};

\node[lab] at (-3,77)  {5. 接口};
\node[layer, text width=150mm] (api) at (78,77) {FastAPI:\ \ \texttt{/search}\ \ \texttt{/chat}\ \ \texttt{/trends}\ \ \texttt{/recommend}\ \ \texttt{/graph}\ \ \texttt{/agent/stream}};

\node[lab] at (-3,94)  {6. 交互};
\node[band, text width=90mm] (fe) at (50,94) {\textbf{终端 AI 助手} \texttt{make tui}(自主工具编排 · SSE 流式)\\{\scriptsize 当前主客户端为 Go TUI;浏览器侧仅保留品牌宣传/下载页边界}};
\node[band, text width=56mm] (llm) at (125,94) {可替换大模型(vLLM/LM Studio)\\{\scriptsize OpenAI 兼容 · 仅生成层}};

\draw[arrow] (raw) -- (proc);
\draw[arrow] (proc) -- (ana);
\draw[arrow] (proc) -- (model);
\draw[arrow] (ana) -- (pg);
\draw[arrow] (model) -- (pg);
\draw[arrow] (pg) -- (api);
\draw[arrow] (api) -- (fe);
\draw[arrow] (api) -- (llm);
\end{tikzpicture}
\caption{SciScope 总体架构:采集、治理、分析/模型、服务、接口与交互逐层衔接;交互层以\textbf{终端智能体}自主编排工具,大模型仅作可替换的生成层。数据来源:系统实现、Makefile 产物与交付目录映射。}
\label{fig:project-architecture}
\end{figure}

\begin{figure}[!htbp]
\centering
\figasset[0.94\textwidth]{system_capability_map.png}
\caption{SciScope 能力地图:从可复现数据资产到证据接地智能体,再到终端产品化客户端。数据来源:系统实现与交付目录映射。}
\label{fig:capability-map}
\end{figure}

\textbf{设计原则}:(1) 离线可复现,每层产物由 \texttt{make} 目标生成;(2) 大模型可替换,生成层走 OpenAI 兼容接口;(3) 证据可追溯,回答必带引用;(4) 本地优先,数据与索引均在本地,支持涉密场景。

\subsection{核心链路:证据接地的检索增强问答}
问答不直接交给大模型臆测,而是\textbf{先检索证据、再据证生成},每条回答都可追溯到具体论文:
见图~\ref{fig:grounded-rag-flow}。

\begin{figure}[!htbp]
\centering
\begin{tikzpicture}[
  x=1mm, y=1mm, font=\scriptsize,
  b/.style={draw=SciLine, rounded corners=1.2pt, line width=0.55pt, align=center, inner sep=3.2pt, minimum height=9mm, fill=white, text=SciInk},
  bb/.style={b, fill=SciPanel},
  arrow/.style={-{Latex[length=1.8mm]}, line width=0.6pt, draw=SciMuted}
]
\node[b, text width=24mm] (q) at (0,0) {用户提问\\(中/英文)};
\node[b, text width=22mm] (fts) at (40,8) {字面检索\\FTS / tsvector};
\node[b, text width=22mm] (vec) at (40,-8) {语义检索\\pgvector 向量};
\node[bb, text width=22mm] (rrf) at (78,0) {RRF 融合\\去重到论文};
\node[bb, text width=22mm] (ev) at (114,0) {证据卡\\标题/年份/片段};
\node[b, text width=24mm] (llm) at (151,0) {大模型据证生成\\\textbf{带引用回答}};
\draw[arrow] (q) -- (fts); \draw[arrow] (q) -- (vec);
\draw[arrow] (fts) -- (rrf); \draw[arrow] (vec) -- (rrf);
\draw[arrow] (rrf) -- (ev); \draw[arrow] (ev) -- (llm);
\end{tikzpicture}
\caption{混合检索 + 证据接地生成链路(中文问句可跨语言命中英文文献)}
\label{fig:grounded-rag-flow}
\end{figure}

\textbf{检索融合口径}: 字面检索与语义检索分别给出排序后, 系统使用倒数排名融合(Reciprocal Rank Fusion, RRF)合并多路结果:
\formulaline{
\operatorname{RRF}(d)=\sum_{r\in R}\frac{1}{k+\operatorname{rank}_r(d)},\quad k=60
}{其中 $R$ 表示参与融合的检索路由, $d$ 为候选论文或片段;该口径能让字面命中和语义相关共同进入证据池。}

\subsection{智能体编排层:从"固定管线"到"自主 agent"}
上节的检索增强生成(Retrieval-Augmented Generation, RAG)问答是\textbf{固定流程}。在其之上,我们构建了\textbf{自主科研智能体}:大模型不再被动接收检索结果,而是\textbf{自己规划步骤、自主选择并调用工具、观察结果后决定下一步},必要时\textbf{自我纠错重试}。图~\ref{fig:agent-loop} 与图~\ref{fig:agent-trace} 展示了从后端编排到终端可视化的同一条事件链。

\begin{figure}[!htbp]
\centering
\begin{tikzpicture}[
  x=1mm, y=1mm, font=\scriptsize,
  b/.style={draw=SciLine, rounded corners=1.2pt, line width=0.55pt, align=center, inner sep=3pt, minimum height=9mm, fill=white, text=SciInk},
  bb/.style={b, fill=SciPanel},
  arrow/.style={-{Latex[length=1.8mm]}, line width=0.6pt, draw=SciMuted},
  loop/.style={-{Latex[length=1.8mm]}, line width=0.6pt, draw=SciTealDark, dashed}
]
\node[b, text width=18mm] (q) at (0,0) {用户问题};
\node[bb, text width=18mm] (plan) at (28,0) {1. 规划\\plan};
\node[b, text width=20mm] (sel) at (58,0) {2. 大模型\\选择工具};
\node[b, text width=20mm] (exec) at (90,0) {3. 并行执行\\只读工具};
\node[b, text width=18mm] (obs) at (121,0) {4. 观察\\结果};
\node[bb, text width=20mm] (ans) at (151,0) {6. 据证作答\\带引用};
\draw[arrow] (q) -- (plan); \draw[arrow] (plan) -- (sel);
\draw[arrow] (sel) -- (exec); \draw[arrow] (exec) -- (obs);
\draw[loop] (obs) to[out=110,in=70] node[above,font=\tiny,text=SciTealDark]{需更多证据则继续} (sel);
\draw[arrow] (obs) -- (ans);
\draw[loop] (ans) to[out=-120,in=-60] node[below,font=\tiny,text=SciTealDark]{5. 反思自评:证据不足->重写检索重试} (sel);
\end{tikzpicture}
\caption{自主智能体循环:规划 → 选工具 → 并行执行 → 观察 →(循环)→ 反思自评 → 据证作答}
\label{fig:agent-loop}
\end{figure}

\begin{figure}[!htbp]
\centering
\figasset[0.94\textwidth]{agent_trace_timeline.png}
\caption{TUI 中可见的智能体执行轨迹:科研工作流阶段、工具调用、证据卡、自检与最终回答均通过流式事件呈现,便于评委在不阅读源码的情况下复核系统不是黑盒问答。}
\label{fig:agent-trace}
\end{figure}

\textbf{核心机制}:
\begin{itemize}
\item \textbf{显式规划}:复杂问题先生成可执行研究计划,再据计划逐步执行;计划对用户与评委\textbf{可见},避免"直接给答案"的黑盒感。
\item \textbf{工具自主编排}:大模型按需调用检索、推荐、趋势、图谱、论断核查等 10 类运行时能力,可\textbf{多步链式}(例如先检索得到真实论文,再据此推荐相近论文或分析趋势);只读能力可并行执行,相同参数自动去重防止空转。
\item \textbf{专业技能工作流}:终端命令把论断核查、文献综述、趋势分析和论文推荐展开为可审计流程模板,约束关键任务必须先取证、再综合。
\item \textbf{反思与自评}:回答后由模型自评"是否充分、关键论断是否都有证据支撑",不足则改写检索\textbf{重试一次}。
\item \textbf{上下文压缩}:长会话自动压缩历史工具结果,避免超出上下文窗口。
\item \textbf{会话态与错误恢复}:每轮交互携带会话标识,同一研究线程内可保存上下文并触发错误恢复,便于连续调研。
\item \textbf{服务/客户端分离}:智能体核心以流式接口输出阶段、耗时、会话与重试标记;当前客户端为 Go TUI,通过同一契约消费智能体核心,后续也可替换为其他界面。
\item \textbf{可插拔生成层}:生成大模型经 OpenAI 兼容接口接入,一处配置即可在\textbf{本地 7B(离线/涉密)}与\textbf{DeepSeek 云端(更强质量)}之间切换,检索、接地与编排逻辑保持不变——模型成为可替换组件而非锁定项。
\item \textbf{工具契约 + 执行前校验门}:每类能力声明只读性、参数校验与结果上限;需要论文主键的动作必须使用真实论文编号,\textbf{在访问数据库之前即拦截模型编造的 ID} 并提示其先检索,从机制上消除"凭空捏造主键"的失败模式。显式技能若首轮跳过必需取证步骤,运行时会强制补上主动作;默认工具预算后强制综合,避免空跑和无限检索。
\item \textbf{自适应检索 + 可观测终止}:模型自行判断是否需要检索(概念/常识类问题直接作答而非空转检索),且检索讲求效率;每轮返回\textbf{终止原因与 token 用量},使智能体的停止条件可解释、调用成本可计量。
\end{itemize}

\begin{plainbox}
\textbf{当前产品化路线}:智能体由单一可观测状态图编排,无运行时开关。底层节点被翻译为用户可理解的科研工作流阶段;TUI 在运行中显示阶段条,并用分组时间线回看证据链。会话标识支持同线程恢复与重试。\textbf{MCP 双向互通与专员子智能体已落地}(见后文"开放与可扩展"小节);下一阶段聚焦持久化 checkpoint 与 Human-in-the-Loop。10 类运行时能力通过同一契约维护,确保报告中的能力清单与实际系统不漂移。
\end{plainbox}

\begin{tableblock}
\begin{tabularx}{0.96\textwidth}{>{\bfseries}p{30mm}XX}
\toprule
用户可见阶段 & 系统动作 & TUI 呈现 \\
\midrule
理解问题 & 解析问题、装载会话上下文 & 显示任务入口与问题摘要 \\
制定研究计划 & 拆解为可执行研究步骤 & 展示计划,评委可直接复核 \\
推理与检索决策 & 判断是否需要取证、调用哪类能力 & 显示当前阶段和后续动作 \\
证据检索 & 调用检索、趋势、推荐、核查等能力 & 返回证据卡和来源信息 \\
自检修正 & 判断证据是否充分,必要时重试 & 标明修正原因与边界 \\
综合回答 & 生成"研究结论 + 证据工具 + 过程入口" & 工具日志不混入主回答 \\
\bottomrule
\end{tabularx}
\end{tableblock}

\begin{tableblock}
\begin{tabularx}{0.96\textwidth}{p{38mm}X}
\toprule
\textbf{运行时能力} & \textbf{评委可见价值} \\
\midrule
跨语言文献检索 & 中文或英文问题均可命中英文论文,输出论文、年份、片段与来源 \\
趋势预测 & 给出关键词/主题热度、增长趋势、生命周期与下一年外推,并标注不确定边界 \\
论文推荐 & 基于语义、关键词、作者关系与时近性给出可解释推荐,不是单纯相似度排序 \\
论文详情与对比 & 对真实论文编号取详情或做多维对比,用于综述与立项查新 \\
领域综述素材 & 检索代表论文并组织成综述素材,支撑快速调研 \\
知识图谱查询 & 展示作者合作、关键词共现、主题社区和 ego 子图,支撑潜在合作者识别 \\
引文导出 & 输出规范 BibTeX 条目,服务论文写作与报告归档 \\
论断核查 & 用跨语言语义接地度判定强支持/部分支持/证据不足,并给出可追溯证据 \\
专员子智能体 & 将综述、趋势、批判核查等聚焦任务交给受限工具集的专员完成 \\
\bottomrule
\end{tabularx}
\end{tableblock}

\subsection{开放与可扩展:MCP 双向互通 + 专员子智能体}
在固定的领域工具之外,智能体层做了三项可扩展性升级,使 SciScope 从"封闭工具集"走向\textbf{开放智能体生态}:
\begin{itemize}
\item \textbf{① SciScope 作为 MCP 服务端}:当前能力可经模型上下文协议(Model Context Protocol, MCP)对外暴露,\textbf{任意 MCP 客户端}(Claude Desktop、Cursor 等)可直接调用我们的接地语料检索与论断核查——把护城河(16 万语料)开放成生态入口而非自用孤岛。
\item \textbf{② 消费外部 MCP 工具}:智能体可将外部 MCP 服务端能力并入统一调度。\textbf{新增能力(如联网取全文、Web 检索)由"插入 MCP 服务"获得,而非改造核心系统},核心工具保持精简、聚焦治理语料。
\item \textbf{③ 专员子智能体(多 Agent 分工)}:复杂多面任务可下派给\textbf{综述员 / 趋势分析师 / 批判核查员}三类专员,各携受限能力集独立完成子任务;专员不再继续委派,\textbf{结构上杜绝递归}。
\end{itemize}
\begin{plainbox}
\textbf{设计取舍(据实)}:不堆砌与科研场景无关的通用编程工具;核心保持领域工具精简,长尾能力靠 MCP 与子智能体扩展——与主流生产级智能体"核心工具精、扩展靠 MCP/sub-agent"的工程哲学一致。
\end{plainbox}

\subsection{方案功能}
\begin{itemize}
\item \textbf{科研智能体}:大模型\textbf{自主规划并编排}检索、核查、趋势、推荐、图谱等能力,多步推理、自我纠错、据证作答;终端助手提供流式交互,并把常见科研任务展开为可审计工具链。
\item \textbf{文献问答}:\textbf{GraphRAG}(查询抽取关键词实体 → 沿共现图扩展邻居 → query expansion)+ 混合检索证据 → 本地大模型据证生成带引用回答 → \textbf{答案可被证据支撑性校验}(低置信明确提示)。
\item \textbf{混合检索}:全文检索 + 向量检索 + RRF 融合,去重到论文级,返回命中片段。
\item \textbf{趋势预测}:关键词/主题的增长率、burst、生命周期与下一年线性外推(含不确定区间)。
\item \textbf{论文推荐}:语义相似 + 关键词重叠 + 作者关系 + 时近性融合,输出\textbf{可解释因子}。
\item \textbf{知识图谱}:作者合作图、关键词共现图、论文-主题图;作者支持实时 ego 子图。
\end{itemize}

\subsection{关键技术}
\begin{itemize}
\item \textbf{向量化}:本地 multilingual-e5-base 嵌入器(768 维,中英跨语言),fp16 加速、断点续传;向量入 pgvector。
\item \textbf{混合检索与 RRF}:全文索引 + 余弦向量索引,两路按 $1/(k+\text{rank})$ 融合(k=60)。
\item \textbf{趋势模型}:逐年归一化词频上最小二乘外推 + 残差不确定带;主题模型 LDA/NMF 各 40 主题。趋势热度 $h_{k,y}=n_{k,y}/N_y$,预测采用 $\hat{h}_{k,y+1}=a_k y+b_k$ 作为排序信号,不作确定性预言。
\item \textbf{推荐模型}:论文级平均向量 + 多信号加权,逐条给出可解释因子;重排采用最大边际相关(Maximum Marginal Relevance, MMR),即 $\lambda sim(q,d)-(1-\lambda)\max_{d'\in S}sim(d,d')$,兼顾相关性与多样性。
\item \textbf{知识图谱}:NetworkX 中心度/社区,百万级边按中心度剪枝导出 JSON,作者 ego 图走 SQL 实时查询。
\item \textbf{生成层}:OpenAI 兼容本地服务,输出限长防失控,可平滑切换更大/远程模型。
\item \textbf{智能体编排}:基于 LangGraph StateGraph 的科研智能体运行时。理解问题、制定计划、选择能力、执行工具、观察结果、自检修正、综合回答被拆为可观察节点;底层使用本地 Qwen2.5-7B,支持流式选工具、只读工具并行执行、调用去重、规划/反思自校验、会话 checkpoint 与同线程重试;事件流向用户呈现科研工作流阶段。
\end{itemize}

\textbf{证据接地口径}:论断核查不把语义相似度包装为形式逻辑证明, 而是将论断 $c$ 与证据片段 $e$ 的接地强度表达为可解释评分:
\formulaline{
g(c,e)=\alpha\,sim(c,e)+\beta\,source(e)+\gamma\,recency(e)
}{该评分用于支持等级与证据排序;最终结论仍必须展示出处,并由 reflect 节点检查关键论断是否有证据支撑。}

\subsection{实现过程}
\textbf{数据生成与采集} → 6 源按领域/年份抓取,形成可追溯的原始分区底账。\\
\textbf{数据组织与治理} → 规范化字段;标题优先去重(合并跨源/预印本同篇,留最全一条)+ front-matter 剔除 + PDF 垃圾/ NUL 清洗;168{,}447 条分析语料进一步形成 159{,}164 篇处理后语料文件,运行时库表当前装载 159{,}187 篇论文。\\
\textbf{信息汲取与建模} → 切片入库,计算片段/论文向量,拟合趋势/推荐/图谱/主题;关键词统一过滤 arXiv 码与超泛标签。模型资产可由一键流水线重建。\\
\textbf{服务与交互} → FastAPI 暴露检索/问答/趋势/推荐/图谱及\textbf{智能体流式}接口;终端 AI 助手自主编排工具;数据库不可用时自动回退内存检索,保证演示鲁棒。

\subsection{交付结构与复现索引}
从评审视角看,本项目按"数据 → 模型 → 服务/RAG → 智能体 → 接口 → 客户端"\textbf{六层交付};职责单一、依赖自底向上:\textbf{Python 承载全部智能}(从数据到 agent),\textbf{Go 仅作终端客户端}(经流式接口消费,可替换、不绑定)。源码路径不作为正文卖点,仅作为复现索引列示。

\begin{tableblock}
\begin{tabularx}{0.96\textwidth}{>{\bfseries}p{22mm}XX}
\toprule
分层 & 评委可见成果 & 复现索引 \\
\midrule
1. 数据层 & 6 源采集、规范化、标题优先去重、分析资产与数据报告 & \texttt{data\_pipeline/}, \texttt{src/harvest}, \texttt{src/analysis} \\
2. 模型层 & 片段/论文向量、趋势、推荐、图谱等\textbf{模型文件} & \texttt{src/models}, \texttt{models/}, \texttt{output/graphs} \\
3. 服务 / RAG & 混合检索、GraphRAG、证据接地、可解释推荐与趋势查询 & \texttt{backend/app/services} \\
4. 智能体 & 可观测状态图:规划、选能力、取证、观察、自检、据证作答 & \texttt{backend/app/agent}, \texttt{.sciscope/skills} \\
5. 接口 & FastAPI 流式接口与常规检索/问答/趋势/推荐/图谱接口 & \texttt{backend/app/api} \\
6. 客户端 & Go 终端客户端;展示科研阶段条、时间线、会话导出与重试 & \texttt{tui/} \\
\bottomrule
\end{tabularx}
\end{tableblock}
